% ===========================================================================
%  Catastrophic Health Expenditure -- NLSS IV
%  Impoverishment from out-of-pocket health spending
%  Generated: April 2026
% ===========================================================================
%
% Required packages (already in your preamble for the other tables):
%   \usepackage{booktabs}
%
% Drop this file into your Overleaf project and % ===========================================================================
%  Catastrophic Health Expenditure -- NLSS IV
%  Impoverishment from out-of-pocket health spending
%  Generated: April 2026
% ===========================================================================
%
% Required packages (already in your preamble for the other tables):
%   \usepackage{booktabs}
%
% Drop this file into your Overleaf project and % ===========================================================================
%  Catastrophic Health Expenditure -- NLSS IV
%  Impoverishment from out-of-pocket health spending
%  Generated: April 2026
% ===========================================================================
%
% Required packages (already in your preamble for the other tables):
%   \usepackage{booktabs}
%
% Drop this file into your Overleaf project and % ===========================================================================
%  Catastrophic Health Expenditure -- NLSS IV
%  Impoverishment from out-of-pocket health spending
%  Generated: April 2026
% ===========================================================================
%
% Required packages (already in your preamble for the other tables):
%   \usepackage{booktabs}
%
% Drop this file into your Overleaf project and \input{impoverishment_table.tex}
% next to where you discuss the OOP / Pen's Parade results.
% ===========================================================================


\subsection*{Out-of-pocket health spending and impoverishment}

Following the Wagstaff and van Doorslaer (2003) framework, we measure the
impoverishing impact of out-of-pocket (OOP) health spending as the increase
in poverty headcount once households' annualised OOP is netted out of
consumption. We use the official NLSS IV per-capita poverty line of
NPR~72{,}908 per person per year, applied to the deflated per-capita
consumption variable (\texttt{pcep}) that the Central Bureau of Statistics
uses to compute the official poverty headcount; per-capita OOP is deflated
by the same household-level Paasche price index for unit consistency. By
construction, gross consumption against the line reproduces the official
20.27\% poverty headcount; netting OOP out of \texttt{pcep} pushes an
additional 3.13 percentage points of the population below the line, equivalent
to roughly 899{,}000 Nepalis pulled into poverty each year by health-related
out-of-pocket expenditure.

For continuity with the earlier CHE poverty draft in the previous try folder,
we also retain the simple nominal per-capita takeaway as a sensitivity check:
using \texttt{pctot\_consumption} against the same line gives a 29.0\%
poverty headcount before OOP and 31.6\% after OOP, an OOP-induced increase of
2.6 percentage points. This sensitivity is not used as the official
impoverishment estimate because \texttt{pctot\_consumption} is nominal whereas
the NLSS poverty line is applied to the deflated \texttt{pcep} welfare measure.

\begin{table}[htbp]
\centering
\caption{OOP-induced impoverishment, Nepal NLSS IV (2022/23)}
\label{tab:impoverishment}
\small
\begin{tabular}{lr}
\hline\hline
\textbf{Indicator} & \textbf{Value} \\
\hline
Poverty headcount, before OOP (gross \texttt{pcep})         & 20.27\% \\
Poverty headcount, after OOP (\texttt{pcep} net of real OOP) & 23.40\% \\
Impoverishment effect of OOP                                & $+3.13$ pp \\
\addlinespace
Households crossing the poverty line due to OOP             & 264 \\
Population pushed below the poverty line due to OOP         & $\sim$899{,}000 \\
\addlinespace
Previous-try nominal per-capita sensitivity                 & 29.0\% $\rightarrow$ 31.6\% ($+2.6$ pp) \\
\hline\hline
\multicolumn{2}{p{12cm}}{\footnotesize \textit{Notes:} Population estimates use NLSS IV individual sampling weights ($\sum \texttt{ind\_wt} \approx 28.74$ million). Welfare measure is \texttt{pcep} -- per-capita real consumption deflated by the Paasche spatial price index, as used in NLSS's official poverty calculation. Out-of-pocket health spending is communicable-disease/injury cost in the past 30 days (\texttt{hh\_comm\_total\_30d}) annualised by $\times 12$, divided by household size, then deflated by the same Paasche index. Poverty line: NPR~72{,}908/yr per capita (\texttt{pline}). Source: NLSS IV.} \\
\end{tabular}
\end{table}

% next to where you discuss the OOP / Pen's Parade results.
% ===========================================================================


\subsection*{Out-of-pocket health spending and impoverishment}

Following the Wagstaff and van Doorslaer (2003) framework, we measure the
impoverishing impact of out-of-pocket (OOP) health spending as the increase
in poverty headcount once households' annualised OOP is netted out of
consumption. We use the official NLSS IV per-capita poverty line of
NPR~72{,}908 per person per year, applied to the deflated per-capita
consumption variable (\texttt{pcep}) that the Central Bureau of Statistics
uses to compute the official poverty headcount; per-capita OOP is deflated
by the same household-level Paasche price index for unit consistency. By
construction, gross consumption against the line reproduces the official
20.27\% poverty headcount; netting OOP out of \texttt{pcep} pushes an
additional 3.13 percentage points of the population below the line, equivalent
to roughly 899{,}000 Nepalis pulled into poverty each year by health-related
out-of-pocket expenditure.

For continuity with the earlier CHE poverty draft in the previous try folder,
we also retain the simple nominal per-capita takeaway as a sensitivity check:
using \texttt{pctot\_consumption} against the same line gives a 29.0\%
poverty headcount before OOP and 31.6\% after OOP, an OOP-induced increase of
2.6 percentage points. This sensitivity is not used as the official
impoverishment estimate because \texttt{pctot\_consumption} is nominal whereas
the NLSS poverty line is applied to the deflated \texttt{pcep} welfare measure.

\begin{table}[htbp]
\centering
\caption{OOP-induced impoverishment, Nepal NLSS IV (2022/23)}
\label{tab:impoverishment}
\small
\begin{tabular}{lr}
\hline\hline
\textbf{Indicator} & \textbf{Value} \\
\hline
Poverty headcount, before OOP (gross \texttt{pcep})         & 20.27\% \\
Poverty headcount, after OOP (\texttt{pcep} net of real OOP) & 23.40\% \\
Impoverishment effect of OOP                                & $+3.13$ pp \\
\addlinespace
Households crossing the poverty line due to OOP             & 264 \\
Population pushed below the poverty line due to OOP         & $\sim$899{,}000 \\
\addlinespace
Previous-try nominal per-capita sensitivity                 & 29.0\% $\rightarrow$ 31.6\% ($+2.6$ pp) \\
\hline\hline
\multicolumn{2}{p{12cm}}{\footnotesize \textit{Notes:} Population estimates use NLSS IV individual sampling weights ($\sum \texttt{ind\_wt} \approx 28.74$ million). Welfare measure is \texttt{pcep} -- per-capita real consumption deflated by the Paasche spatial price index, as used in NLSS's official poverty calculation. Out-of-pocket health spending is communicable-disease/injury cost in the past 30 days (\texttt{hh\_comm\_total\_30d}) annualised by $\times 12$, divided by household size, then deflated by the same Paasche index. Poverty line: NPR~72{,}908/yr per capita (\texttt{pline}). Source: NLSS IV.} \\
\end{tabular}
\end{table}

% next to where you discuss the OOP / Pen's Parade results.
% ===========================================================================


\subsection*{Out-of-pocket health spending and impoverishment}

Following the Wagstaff and van Doorslaer (2003) framework, we measure the
impoverishing impact of out-of-pocket (OOP) health spending as the increase
in poverty headcount once households' annualised OOP is netted out of
consumption. We use the official NLSS IV per-capita poverty line of
NPR~72{,}908 per person per year, applied to the deflated per-capita
consumption variable (\texttt{pcep}) that the Central Bureau of Statistics
uses to compute the official poverty headcount; per-capita OOP is deflated
by the same household-level Paasche price index for unit consistency. By
construction, gross consumption against the line reproduces the official
20.27\% poverty headcount; netting OOP out of \texttt{pcep} pushes an
additional 3.13 percentage points of the population below the line, equivalent
to roughly 899{,}000 Nepalis pulled into poverty each year by health-related
out-of-pocket expenditure.

For continuity with the earlier CHE poverty draft in the previous try folder,
we also retain the simple nominal per-capita takeaway as a sensitivity check:
using \texttt{pctot\_consumption} against the same line gives a 29.0\%
poverty headcount before OOP and 31.6\% after OOP, an OOP-induced increase of
2.6 percentage points. This sensitivity is not used as the official
impoverishment estimate because \texttt{pctot\_consumption} is nominal whereas
the NLSS poverty line is applied to the deflated \texttt{pcep} welfare measure.

\begin{table}[htbp]
\centering
\caption{OOP-induced impoverishment, Nepal NLSS IV (2022/23)}
\label{tab:impoverishment}
\small
\begin{tabular}{lr}
\hline\hline
\textbf{Indicator} & \textbf{Value} \\
\hline
Poverty headcount, before OOP (gross \texttt{pcep})         & 20.27\% \\
Poverty headcount, after OOP (\texttt{pcep} net of real OOP) & 23.40\% \\
Impoverishment effect of OOP                                & $+3.13$ pp \\
\addlinespace
Households crossing the poverty line due to OOP             & 264 \\
Population pushed below the poverty line due to OOP         & $\sim$899{,}000 \\
\addlinespace
Previous-try nominal per-capita sensitivity                 & 29.0\% $\rightarrow$ 31.6\% ($+2.6$ pp) \\
\hline\hline
\multicolumn{2}{p{12cm}}{\footnotesize \textit{Notes:} Population estimates use NLSS IV individual sampling weights ($\sum \texttt{ind\_wt} \approx 28.74$ million). Welfare measure is \texttt{pcep} -- per-capita real consumption deflated by the Paasche spatial price index, as used in NLSS's official poverty calculation. Out-of-pocket health spending is communicable-disease/injury cost in the past 30 days (\texttt{hh\_comm\_total\_30d}) annualised by $\times 12$, divided by household size, then deflated by the same Paasche index. Poverty line: NPR~72{,}908/yr per capita (\texttt{pline}). Source: NLSS IV.} \\
\end{tabular}
\end{table}

% next to where you discuss the OOP / Pen's Parade results.
% ===========================================================================


\subsection*{Out-of-pocket health spending and impoverishment}

Following the Wagstaff and van Doorslaer (2003) framework, we measure the
impoverishing impact of out-of-pocket (OOP) health spending as the increase
in poverty headcount once households' annualised OOP is netted out of
consumption. We use the official NLSS IV per-capita poverty line of
NPR~72{,}908 per person per year, applied to the deflated per-capita
consumption variable (\texttt{pcep}) that the Central Bureau of Statistics
uses to compute the official poverty headcount; per-capita OOP is deflated
by the same household-level Paasche price index for unit consistency. By
construction, gross consumption against the line reproduces the official
20.27\% poverty headcount; netting OOP out of \texttt{pcep} pushes an
additional 3.13 percentage points of the population below the line, equivalent
to roughly 899{,}000 Nepalis pulled into poverty each year by health-related
out-of-pocket expenditure.

For continuity with the earlier CHE poverty draft in the previous try folder,
we also retain the simple nominal per-capita takeaway as a sensitivity check:
using \texttt{pctot\_consumption} against the same line gives a 29.0\%
poverty headcount before OOP and 31.6\% after OOP, an OOP-induced increase of
2.6 percentage points. This sensitivity is not used as the official
impoverishment estimate because \texttt{pctot\_consumption} is nominal whereas
the NLSS poverty line is applied to the deflated \texttt{pcep} welfare measure.

\begin{table}[htbp]
\centering
\caption{OOP-induced impoverishment, Nepal NLSS IV (2022/23)}
\label{tab:impoverishment}
\small
\begin{tabular}{lr}
\hline\hline
\textbf{Indicator} & \textbf{Value} \\
\hline
Poverty headcount, before OOP (gross \texttt{pcep})         & 20.27\% \\
Poverty headcount, after OOP (\texttt{pcep} net of real OOP) & 23.40\% \\
Impoverishment effect of OOP                                & $+3.13$ pp \\
\addlinespace
Households crossing the poverty line due to OOP             & 264 \\
Population pushed below the poverty line due to OOP         & $\sim$899{,}000 \\
\addlinespace
Previous-try nominal per-capita sensitivity                 & 29.0\% $\rightarrow$ 31.6\% ($+2.6$ pp) \\
\hline\hline
\multicolumn{2}{p{12cm}}{\footnotesize \textit{Notes:} Population estimates use NLSS IV individual sampling weights ($\sum \texttt{ind\_wt} \approx 28.74$ million). Welfare measure is \texttt{pcep} -- per-capita real consumption deflated by the Paasche spatial price index, as used in NLSS's official poverty calculation. Out-of-pocket health spending is communicable-disease/injury cost in the past 30 days (\texttt{hh\_comm\_total\_30d}) annualised by $\times 12$, divided by household size, then deflated by the same Paasche index. Poverty line: NPR~72{,}908/yr per capita (\texttt{pline}). Source: NLSS IV.} \\
\end{tabular}
\end{table}
